\documentclass[11pt,a4paper]{article}

\usepackage[margin=2.3cm]{geometry}
\usepackage{graphicx}
\usepackage{booktabs}
\usepackage{amsmath}
\usepackage{amssymb}
\usepackage{caption}
\usepackage{subcaption}
\usepackage{float}
\usepackage{array}
\usepackage{enumitem}
\usepackage{xcolor}
\usepackage{hyperref}
\usepackage{url}

\graphicspath{{figures/}}
\hypersetup{colorlinks=true,linkcolor=blue!55!black,urlcolor=blue!55!black,citecolor=blue!55!black}
\captionsetup{font=small,labelfont=bf}
\setlength{\parskip}{0.35em}
\setlength{\parindent}{0pt}
\setlength{\emergencystretch}{2em}

\title{\textbf{Thermal Object Detection in Adverse Weather}\\[0.3em]
\large A Consolidated Report on Real-Weather Baselines, Synthetic-Weather Generation,\\
and Learned Image Enhancement}
\author{GATE-IR Project}
\date{15 November 2026}

\begin{document}
\maketitle

\begin{abstract}
This report consolidates the experiments carried out in the GATE-IR project on
thermal (infrared) object detection under adverse weather. We study three
datasets in increasing order of control: \textbf{UNIRI-TID}, which contains real
clear, fog and rain sequences; \textbf{M3FD}, a paired infrared/visible dataset
for which fog and rain were generated synthetically by three different methods;
and \textbf{FLIR ADAS v2}, a large thermal dataset with no real fog, for which
fog was generated with CycleGAN and with a diffusion-based translation model.
Across these datasets we compare detectors, weather-synthesis methods, classical
image filters, a learned U-Net dehazer, and a learned clear/fog gate.

The main findings are: (i) on real weather, the clear-to-fog accuracy gap is
small once training and evaluation use the same frames, and mixed-weather
training largely closes it; (ii) on M3FD, deterministic synthetic fog is essentially free after mixed
training, whereas learned (CycleGAN) fog retains a residual gap of about 0.05; (iii) on FLIR, only the learned
U-Net dehazer helps, classical filters do not, and a 100-epoch matched comparison
improves fog mAP@50 from 0.2430 to 0.2856 (+17.5\%) while reducing the
clear-to-fog gap by 43\%; (iv) the dehazer does not transfer across generators,
and even a dehazer retrained on the target generator remains below the
unfiltered baseline; and (v) gating is an efficiency technique, not a reliable
accuracy technique: in an indicative 40-epoch comparison always-on filtering
scored higher, but that comparison used different detector sizes (YOLOv8s versus
YOLOv12s) and is therefore not controlled; and (vi) an eight-condition extension
on FLIR shows that snow, sensor degradation, defocus and static obstruction
cause clear, measurable gaps, and that mixed training closes 52--78\,\% of them.
Small-target pixelation is the exception: it is a strong learnable cue, producing
a negative gap rather than a robustness gain.

All numbers in this report were extracted from result files preserved with the
project. No experiment was re-run while the report was written.
\end{abstract}

\tableofcontents
\newpage

%======================================================================
\section{Introduction and scope}

Thermal cameras measure emitted infrared radiation and therefore see at night
and in low light, which makes them attractive for safety and automotive
perception. Their weakness is adverse weather: fog, rain and snow scatter or
absorb radiation, reduce contrast, and erase the fine edges that object
detectors rely on. The GATE-IR project asks a practical question:

\begin{quote}
\emph{Can a cheap, learned preprocessing stage make a thermal detector robust to
adverse weather, and is it better to apply that stage to every frame or only to
frames that a small gate classifies as adverse?}
\end{quote}

We study the question in three settings, reported in this order:

\begin{enumerate}[leftmargin=1.4em]
  \item \textbf{UNIRI-TID} (Section~\ref{sec:uniri}): real clear, fog and rain
        thermal sequences. This is the most realistic setting and establishes
        what the weather gap looks like without synthetic data.
  \item \textbf{M3FD} (Section~\ref{sec:m3fd}): paired infrared/visible scenes
        with synthetically generated fog and rain. This lets us compare three
        very different synthesis methods under a controlled protocol.
  \item \textbf{FLIR ADAS v2} (Section~\ref{sec:flir}): a large thermal driving
        dataset with no real fog. Fog is generated with two learned translation
        models, and the detector, filter and gate are compared on exactly the
        same evaluation frames.
\end{enumerate}

Section~\ref{sec:flir-adverse} extends the FLIR line with eight additional
adverse conditions. It screens a clear-trained detector on each condition and
retrains only on the conditions that create a measurable gap, rather than
assuming that every simulated degradation is equally harmful.

Section~\ref{sec:discussion} discusses the cross-cutting lessons,
Section~\ref{sec:limitations} lists the limitations, and
Section~\ref{sec:repro} explains how the work can be reproduced.

%======================================================================
\section{Terminology}

The report deliberately avoids internal experiment codes. Table~\ref{tab:terms}
defines every term used later.

\begin{table}[H]
\centering
\small
\caption{Terminology used in this report.}
\label{tab:terms}
\begin{tabular}{@{}p{0.26\linewidth}p{0.68\linewidth}@{}}
\toprule
Term & Meaning in this report \\
\midrule
Thermal / infrared image & A single-channel image of emitted heat, usually displayed in grayscale or a false-colour palette. \\
Visible (RGB) image & An ordinary three-channel photograph. M3FD contains paired thermal and visible images. \\
Adverse weather & Fog, rain or snow. In this project, ``adverse'' always means one of these three conditions. \\
Synthetic weather & An image generated from a clear image by a weather-simulation method, rather than captured in real weather. \\
Clear-only training & Training a detector only on clear images, then testing it on adverse images. This measures the raw weather gap. \\
Mixed-weather training & Training on clear images plus adverse images, then testing on adverse images. This is the realistic deployment protocol. \\
Domain shift & Accuracy loss caused by a difference between training and test image statistics (for example, different cameras, scenes or generators), not by weather itself. \\
mAP@50 & Mean average precision with an intersection-over-union threshold of 0.50. \\
mAP@50--95 & Mean average precision averaged over thresholds 0.50, 0.55, \ldots, 0.95. The stricter and more informative detection metric. \\
APm & Average precision for medium-sized objects (COCO definition, 32$^2$ to 96$^2$ pixels). Used on UNIRI-TID. \\
Gap & Clear-image score minus adverse-image score, in the same metric. A smaller gap is better. \\
CycleGAN & An unpaired image-to-image translation model that learns a mapping between two image domains without paired examples. \\
CycleDiff & A diffusion-based translation model that combines a latent diffusion model with cycle consistency. \\
Dehazer & A learned image-to-image network (here, a small U-Net) trained to map an adverse image back towards a clear image. \\
Gamma correction & A point-wise intensity transform that brightens or darkens an image. \\
CLAHE & Contrast-limited adaptive histogram equalisation, a local contrast enhancement filter. \\
DCP & Dark-channel prior, a classical physics-inspired dehazing method. \\
Gate & A small classifier that decides whether a frame is clear or adverse. \\
Gated inference & At test time, the gate sends clear frames directly to the detector and adverse frames through the filter first. \\
Always-on filtering & Every frame, clear or adverse, passes through the filter before detection. \\
Matched-frame evaluation & Scoring clear and adverse conditions on exactly the same frame indices, so that the measured difference is weather, not a different test set. \\
Early stopping & Ending training when the validation score has not improved for a fixed number of epochs (25 here). \\
\bottomrule
\end{tabular}
\end{table}

%======================================================================
\section{Datasets}

\subsection{UNIRI-TID (real weather)}
UNIRI-TID\cite{uniri} is a thermal dataset with real clear, fog and rain sequences. We use a
sequence-disjoint split (no sequence appears in more than one split) and report
results on the test split: \textbf{1,234 images} (535 clear, 195 fog, 504 rain).
The images are false-colour thermal frames; the palette is part of the recorded
data and is identical across conditions. The dataset is small and the fog test
set in particular is limited to two sequences, so absolute fog numbers are less
stable than the clear and rain numbers.

\subsection{M3FD (synthetic weather on paired thermal/visible data)}
M3FD\cite{m3fd} contains paired thermal and visible images of driving scenes with six
object classes. We use a chunk-disjoint temporal split with hash
\texttt{effb9b3090ad2fcb}. Fog and rain are generated synthetically by three
methods (Section~\ref{sec:m3fd-methods}). The clear baseline is evaluated on the
same test frames as the synthetic conditions.

\subsection{FLIR ADAS v2 (synthetic fog on thermal data)}
FLIR ADAS v2\cite{flir} is a large thermal driving dataset with 15 object classes. It has
no real fog. We therefore use the \textbf{3,749-frame FLIR test set} as the
common evaluation set and generate fog with two learned translation models. Every condition (clear, fog, filtered fog) is scored on the
same 3,749 frames with the same labels. This matched-frame protocol is essential:
older runs compared 1,200 in-domain clear frames with 3,749 out-of-domain fog
frames, which overstates the weather gap (Section~\ref{sec:flir-protocol}).


%======================================================================
\section{Methods}

\subsection{Detectors}
All detectors are one-stage YOLO-family models\cite{yolo} trained with the same recipe:
640$\times$640 input, AdamW with automatic learning rate, cosine schedule,
early stopping after 25 epochs without improvement, and a fixed random seed.
Unless stated otherwise, training uses a batch size of 16 on FLIR and 64 on
M3FD, and evaluation uses a confidence threshold of 0.001 and an IoU threshold
of 0.5 for non-maximum suppression. We report mAP@50 and mAP@50--95; on
UNIRI-TID we additionally report APm because small and medium objects dominate
the scene.

\paragraph{Reading the tables.}
All detection scores are fractions in $[0,1]$ and higher is better, except the
gap, which is always \emph{clear minus adverse} and lower is better. mAP@50 and
mAP@50--95 are defined in Table~\ref{tab:terms}; APm is medium-object average
precision. \(n\) is the number of test images, and ``--'' means not applicable.

\subsection{Weather synthesis}
\label{sec:m3fd-methods}
Three families were compared on M3FD:
\begin{description}[leftmargin=1.4em,style=nextline]
  \item[Classical overlays.] A uniform veil, rain streaks or snow particles are
        composited onto the clear thermal image. The method is deterministic,
        fast (minutes), and its severity is controlled explicitly.
  \item[Depth-aware fog.] A Koschmieder fog model is applied with a
        transmission map derived from a depth estimate. This produces spatially
        varying fog that respects scene geometry.
  \item[Learned translation.] CycleGAN\cite{cyclegan} learns an unpaired mapping from clear to
        adverse using real adverse thermal images (DSERT-RoLL) as the target
        domain. On FLIR we also use CycleDiff\cite{cyclediff}, a diffusion-based translator, and
        a 75\,\% blend of clear and translated fog.
\end{description}

\subsection{Enhancement filters}
We compare four ways to make an adverse image easier to detect:
\begin{itemize}[leftmargin=1.4em]
  \item \textbf{Gamma correction:} an adaptive point-wise brightening based on
        the image mean.
  \item \textbf{CLAHE:} contrast-limited adaptive histogram equalisation on the
        lightness channel\cite{clahe}.
  \item \textbf{Dark-channel prior (DCP):} a classical physics-inspired dehazing
        operator\cite{dcp} with a guided-filter refinement.
  \item \textbf{Learned U-Net dehazer:} a small encoder--decoder network trained
        on pairs of (clear, synthetic adverse) images to reconstruct the clear
        image.
\end{itemize}

\subsection{Gating}
A binary clear/adverse classifier is trained on five simple image features
(variance, minimum, maximum, entropy and Laplacian variance) followed by a small
multilayer perceptron. At inference, the gate can route clear frames directly to
the detector and adverse frames through the filter (gated inference), or the
filter can be applied to every frame (always-on filtering).

%======================================================================
\section{UNIRI-TID: real-weather baselines}
\label{sec:uniri}

\subsection{Question}
Before building any synthetic-weather pipeline, we asked how large the weather
gap is when the weather is real. If the gap is already small, a gate or a
heavy restoration network has little accuracy headroom.


\subsection{Detector comparison}
Table~\ref{tab:uniri} compares six detector configurations and one invalid
RT-DETR run; Figure~\ref{fig:uniri} shows representative real-weather frames. All models use the same sequence-disjoint split.

\begin{table}[H]
\centering
\small
\caption{UNIRI-TID test results. Each row is one detector. Columns: overall test
mAP@50--95; APm for clear, fog and rain; and the clear-minus-fog APm gap. Higher
is better except the gap, which is better when smaller.}
\label{tab:uniri}
\begin{tabular}{@{}lccccc@{}}
\toprule
Detector & mAP@50--95 & APm clear & APm fog & APm rain & Clear$-$fog \\
\midrule
YOLOv8s                       & 0.6204 & 0.6623 & 0.5344 & 0.5951 & 0.128 \\
YOLO11s                       & 0.6194 & 0.6585 & 0.5264 & 0.6038 & 0.132 \\
\textbf{YOLOv12s}             & \textbf{0.6296} & \textbf{0.6670} & \textbf{0.5733} & 0.6025 & \textbf{0.094} \\
YOLOv8s + weather augmentation& 0.6278 & 0.6652 & 0.5691 & 0.6035 & 0.096 \\
YOLOv10s                      & 0.6063 & 0.6459 & 0.5316 & 0.5873 & 0.114 \\
YOLOv8s-P2                    & 0.6071 & 0.6480 & 0.4954 & 0.5909 & 0.153 \\
RT-DETR-l (invalid)           & 0.2778 & -- & -- & -- & -- \\
\bottomrule
\end{tabular}
\end{table}

\begin{figure}[H]
\centering
\includegraphics[width=\linewidth]{fig_uniri_conditions.png}
\caption{Real thermal conditions in UNIRI-TID. Rows are clear, fog and rain;
columns are three non-adjacent frames from different sequences (frame~1354 /
sequence~3, frame~2464 / sequence~17, frame~2881 / sequence~19 for clear; frames
4440, 4555 and 4634 for fog; frames 4749, 5403 and 7109 for rain). The images are
false-colour thermal frames; the palette and the on-screen temperature scale are
part of the recorded data.}
\label{fig:uniri}
\end{figure}


\subsection{Key points}
\begin{itemize}[leftmargin=1.4em]
  \item The best detector is YOLOv12s (mAP@50--95 0.6296), and it also has the
        smallest clear-to-fog APm gap (0.094).
  \item Weather augmentation improves fog APm from 0.5344 to 0.5691 at almost no
        cost in clear APm. Mixed-weather training is therefore effective on real
        weather.
  \item mAP@50 is saturated (roughly 0.96--0.97 for all models) and is not a
        useful discriminator; the report uses mAP@50--95 and APm.
  \item The RT-DETR-l run stopped after 29 epochs because of a wall-time limit
        and is reported only to document that it is invalid.
\end{itemize}

A gated-preprocessing variant study was scaffolded on UNIRI-TID, but it did not
produce a completed result table (only logs), so it is not used for conclusions.
The real-weather evidence pointed in a clear direction: because mixed training
already closes most of the gap, the project moved to controlled synthetic weather
where the effect of a filter could be isolated.

%======================================================================
\section{M3FD: comparing weather-synthesis methods}
\label{sec:m3fd}

\subsection{Question}
If synthetic weather is used for training, does the choice of synthesis method
matter? We compare classical overlays, depth-aware fog and CycleGAN under two
protocols:
\begin{itemize}[leftmargin=1.4em]
  \item \textbf{Clear-only training:} train on clear thermal images, test on
        synthetic weather. This measures how far the synthetic weather moves the
        image away from the training distribution.
  \item \textbf{Mixed-weather training:} train on clear plus synthetic weather,
        test on synthetic weather. This measures whether the detector can learn
        the synthetic weather.
\end{itemize}

\begin{figure}[H]
\centering
\includegraphics[width=\linewidth]{fig_m3fd_synthesis.png}
\caption{M3FD synthetic-weather comparison for three non-adjacent frames
(00001, 00060 and 01170). Columns: the original clear thermal image, classical
fog, depth-aware fog, and learned fog from CycleGAN. The clear image is the
reference; the three synthetic columns are generated from it. The classical and
depth methods produce a smooth veil, while CycleGAN produces a darker, more
textured fog.}
\label{fig:m3fd}
\end{figure}

\subsection{Results}
Figure~\ref{fig:m3fd} shows the three synthesis methods on non-adjacent frames,
and Table~\ref{tab:m3fd} reports the mAP of YOLOv12s on the clean split. The clear
baseline is mAP@50 0.268 and mAP@50--95 0.148.

\begin{table}[H]
\centering
\footnotesize
\caption{M3FD synthetic-weather results (YOLOv12s, clean split
\texttt{effb9b3090ad2fcb}). Each row is one (synthesis method, training protocol,
test condition) combination. ``Clear-only'' means trained on clear images and
tested on the condition; ``mixed'' means trained on clear plus the condition.
Columns: mAP@50, mAP@50--95, and the clear-minus-condition gap in mAP@50--95;
lower gap is better.}
\label{tab:m3fd}
\begin{tabular}{@{}llcccc@{}}
\toprule
Method & Training & Condition & mAP@50 & mAP@50--95 & Gap \\
\midrule
Classical   & clear-only & fog light   & 0.268 & 0.148 & 0.001 \\
Classical   & clear-only & fog medium  & 0.243 & 0.135 & 0.013 \\
Classical   & clear-only & fog heavy   & 0.206 & 0.117 & 0.032 \\
Classical   & clear-only & rain light  & 0.261 & 0.144 & 0.005 \\
Classical   & clear-only & rain medium & 0.189 & 0.094 & 0.054 \\
Classical   & clear-only & rain heavy  & 0.079 & 0.041 & 0.108 \\
\midrule
Depth       & clear-only & fog medium  & 0.199 & 0.114 & 0.035 \\
Depth       & clear-only & fog heavy   & 0.097 & 0.072 & 0.076 \\
Depth       & clear-only & rain medium & 0.094 & 0.054 & 0.054 \\
Depth       & clear-only & rain heavy  & 0.011 & 0.011 & 0.137 \\
Depth       & mixed      & clear test  & 0.244 & 0.141 & -- \\
Depth       & mixed      & fog heavy   & 0.237 & 0.140 & $-0.001$ \\
\midrule
CycleGAN    & clear-only & fog         & 0.075 & 0.048 & 0.101 \\
CycleGAN    & clear-only & rain        & 0.185 & 0.105 & 0.044 \\
CycleGAN    & mixed      & clear test  & 0.260 & 0.147 & -- \\
CycleGAN    & mixed      & fog         & 0.193 & 0.098 & 0.049 \\
CycleGAN    & mixed      & rain        & 0.197 & 0.104 & 0.011 \\
\bottomrule\end{tabular}
\end{table}

\subsection{Key points}
\begin{itemize}[leftmargin=1.4em]
  \item \textbf{Deterministic fog is learnable.} Depth-aware fog has a
        clear-only gap of 0.076 but a mixed-training gap of $-0.001$: once the
        detector sees the same deterministic transform during training, the
        weather costs nothing.
  \item \textbf{CycleGAN fog has the largest clear-only gap} (0.101), because its
        texture and contrast differ most from clear images. Mixed training
        reduces that gap to 0.049, but it does not eliminate it.
  \item \textbf{Learned fog retains a residual gap.} Under mixed training the
        remaining gaps are 0.049 for CycleGAN fog, 0.011 for CycleGAN rain and
        $-0.001$ for depth-aware fog. The practical lesson is that the choice of
        synthesis method matters more than the nominal severity: a deterministic
        veil is easier to learn than a learned translation.
  \item \textbf{Rain is harder than fog} for the classical and depth methods:
        heavy rain drops the clear-only mAP@50--95 to 0.041 and 0.011
        respectively. This is consistent with rain streaks destroying thin
        structures rather than merely lowering contrast.
\end{itemize}

%======================================================================
\section{FLIR ADAS v2: learned enhancement under synthetic fog}
\label{sec:flir}

\subsection{Question}
FLIR has no real fog, but it has a large, well-labelled thermal test set. We
therefore generate fog with two learned models and ask three questions:
\begin{enumerate}[leftmargin=1.4em]
  \item Does a learned dehazer improve detection on synthetic fog?
  \item Do classical filters help?
  \item Is it better to filter every frame or to gate the filter?
\end{enumerate}

\subsection{Data and matched-frame protocol}
\label{sec:flir-protocol}
The evaluation set is the 3,749-frame test set. For each frame we have:
the original clear image, a CycleGAN fog image, a CycleDiff-75 fog image (a
75\,\% blend of clear and diffusion-generated fog), and a U-Net-dehazed version
of the CycleDiff-75 image. All conditions share the same labels and the same
frame indices.

This matters because the older protocol compared \textbf{1,200 in-domain clear
frames} with \textbf{3,749 out-of-domain fog frames}. Re-scoring the same fog
model on the same 3,749 clear frames gives clear mAP@50--95 0.1669 and fog
0.1445, a gap of only 0.022. The apparent 0.18 gap in the old protocol was mostly
domain shift, not weather.

\begin{figure}[H]
\centering
\includegraphics[width=\linewidth]{fig_flir_synthesis.png}
\caption{FLIR synthetic fog and dehazing for three non-adjacent frames from
three different videos (Video~A: \texttt{4FRnNpmSmwktFJKjg}, frame~000938;
Video~B: \texttt{ZAtDSNuZZjkZFvMAo}, frame~000896; Video~C:
\texttt{SCiKdG3MqZfiE292B}, frame~007909). Columns: original clear thermal
image, CycleGAN fog, CycleDiff-75 fog, CycleDiff-100 fog, and the U-Net dehazed
version of the CycleDiff-75 image. The 100\,\% diffusion translation is almost
featureless; the 75\,\% blend is the usable one. The dehazer restores much of
the contrast and edge structure of the clear image.}
\label{fig:flir-synthesis}
\end{figure}

\subsection{Weather synthesis quality}
The paired clear-to-fog translation was also scored with image-similarity
metrics: PSNR 14.01\,dB, SSIM 0.636 and LPIPS 0.602 over 3,749 pairs. These
numbers measure how far the translation moves away from the clear source; they
are not a realism score. No FID was computed, so we do not claim a distributional
quality result.

\begin{figure}[H]
\centering
\includegraphics[width=\linewidth]{fig_flir_weather.png}
\caption{The same FLIR frame under clear conditions and two CycleGAN weather
translations (fog and rain). It illustrates the weather-synthesis coverage
available for training; the quantitative experiments focus on fog because it is
the condition with a complete detector comparison.}
\label{fig:flir-weather}
\end{figure}

\paragraph{Training time.}
The final CycleDiff translator continuation (100k to 105k steps) ran from
14:17 to 18:41 on 25 October 2026, i.e.\ 4\,h\,24\,min, measured from the
creation and last-write times of \path{experiments/01_flir_weather_translation/method4_cyclediff/logs/train_100k_105k.log} and the matching
event file. The pipeline was trained in stages: the two VAEs on 9--10 October,
the two latent-diffusion models on 12--15 October, and the cycle translator in
several runs between 15 and 25 October. Because the stages waited in the cluster
queue and resumed several times, the 10--25 October calendar span is not a
compute-time figure; the final translator run is the directly comparable
training cost.

Figure~\ref{fig:flir-synthesis} includes the 100\,\% CycleDiff translation. It is
almost featureless in these examples, and the detector trained on it scored
0.000 mAP in the original evaluation; the experiments therefore use the
75\,\% blend.

\subsection{Matched detector baselines}
Table~\ref{tab:flir-original} reports the matched test-set scores (all 3,749
frames). Both rows are matched clear/fog pairs. The older 1,200-frame in-domain
clear evaluations are omitted because they are not comparable to the
3,749-frame fog test.

\begin{table}[H]
\centering
\footnotesize
\setlength{\tabcolsep}{4pt}
\caption{Matched test-set results for the two usable generators. Every score
uses the same 3,749-frame test set. Each row is one trained model; columns are
clear mAP@50, clear mAP@50--95, fog mAP@50, fog mAP@50--95, and the
clear-minus-fog mAP@50--95 gap. Both rows are matched clear/fog pairs. Rows are
not a model-versus-model comparison: compare clear and fog within a row, not
across rows.}
\label{tab:flir-original}
\begin{tabular}{@{}lccccc@{}}
\toprule
Generator & Clear 50 & Clear 50--95 & Fog 50 & Fog 50--95 & Gap \\
\midrule
CycleGAN      & 0.3055 & 0.1669 & 0.2718 & 0.1445 & 0.0224 \\
CycleDiff-75  & 0.2746 & 0.1491 & 0.2747 & 0.1495 & $-0.0004$ \\
\bottomrule
\end{tabular}
\end{table}

\subsection{Filter screen}
A 60-epoch screen compared five filters on the same matched frames with a
YOLOv8s detector (Table~\ref{tab:flir-screen}). Only the learned U-Net dehazer
improves the fog score; the classical filters either leave it unchanged or make
it worse (Figure~\ref{fig:flir-filters}).

\begin{table}[H]
\centering
\small
\caption{FLIR filter screen (YOLOv8s, 60 epochs, matched 3,749-frame test set).
Each row is one filter applied to the same fog frames; columns are clear and fog
mAP@50--95 and their gap. ``None'' is the unfiltered baseline.}
\label{tab:flir-screen}
\begin{tabular}{@{}lccc@{}}
\toprule
Filter & Clear mAP@50--95 & Fog mAP@50--95 & Gap \\
\midrule
None (baseline)      & 0.1646 & 0.1413 & 0.0233 \\
Gamma                & 0.1588 & 0.1369 & 0.0219 \\
CLAHE                & 0.1636 & 0.1335 & 0.0301 \\
Dark-channel prior   & 0.1651 & 0.1394 & 0.0257 \\
\textbf{U-Net dehazer} & \textbf{0.1676} & \textbf{0.1524} & \textbf{0.0152} \\
Gamma + CLAHE        & 0.1610 & 0.1384 & 0.0227 \\
\bottomrule
\end{tabular}
\end{table}

\begin{figure}[H]
\centering
\includegraphics[width=\linewidth]{fig_flir_filters.png}
\caption{Classical filters applied to the same CycleGAN fog frames shown in
Figure~\ref{fig:flir-synthesis}. Columns: synthetic fog, gamma correction,
CLAHE, dark-channel prior, and gamma followed by CLAHE. The filters change the
global tone but do not recover the scene structure lost to the fog; this matches
their failure to improve detection in Table~\ref{tab:flir-screen}.}
\label{fig:flir-filters}
\end{figure}

\subsection{Main 100-epoch comparison}
The main comparison uses YOLOv12s and the full 100-epoch budget (early stopping
after 25 epochs). Two pipelines are compared:
\begin{itemize}[leftmargin=1.4em]
  \item \textbf{Baseline:} train on clear plus unfiltered CycleGAN fog, detect
        directly.
  \item \textbf{Filtered:} train on clear plus U-Net-dehazed fog, detect
        directly. At deployment, this corresponds to always-on filtering; the
        gated variant is discussed below.
\end{itemize}

\begin{table}[H]
\centering
\small
\caption{FLIR main comparison, YOLOv12s, matched 3,749-frame test set. Each row is
one training/inference pipeline. Columns: clear and fog mAP@50, clear and fog
mAP@50--95, and the clear-minus-fog mAP@50--95 gap.}
\label{tab:flir-main}
\begin{tabular}{@{}lccccc@{}}
\toprule
Pipeline & Clear mAP@50 & Clear mAP@50--95 & Fog mAP@50 & Fog mAP@50--95 & Gap \\
\midrule
Baseline & 0.2849 & 0.1525 & 0.2430 & 0.1266 & 0.0259 \\
\textbf{Filtered (U-Net)} & \textbf{0.3063} & \textbf{0.1683} & \textbf{0.2856} & \textbf{0.1535} & \textbf{0.0148} \\
$\Delta$ & $+0.0214$ & $+0.0158$ & $+0.0426$ & $+0.0269$ & $-42.9\%$ \\
\bottomrule
\end{tabular}
\end{table}

Table~\ref{tab:flir-main} reports the main comparison. The filtered pipeline
improves fog mAP@50 by 17.5\,\% and fog mAP@50--95 by
21.3\,\%, and reduces the clear-to-fog gap by 42.9\,\%. It also improves the
\emph{clear} score, so the gain is not a clear-versus-fog trade-off. Both
pipelines early-stopped: the baseline at epoch 55 (best epoch 30) and the
filtered model at epoch 57 (best epoch 32).

\subsection{Does the dehazer transfer across generators?}
A 40-epoch YOLOv8s study trained the U-Net on CycleGAN fog and tested it on
CycleDiff-75 fog. It hurts: fog mAP@50--95 falls from 0.1619 to 0.1511 and clear
from 0.1664 to 0.1593. Retraining the U-Net on CycleDiff-75 pairs improves the
result (clear 0.1638, fog 0.1560) but still leaves it below the unfiltered
baseline (0.1664 / 0.1619). The reconstruction objective is therefore not
aligned with detection, even when the generator matches.

\subsection{Always-on filtering versus gated inference}
The gate was evaluated as a routing mechanism. On the generator it was trained
on (CycleGAN) it separates fog from clear almost perfectly; on the other
generator (CycleDiff-75) it misses 33.9\,\% of fog frames (Table~\ref{tab:gate}).

\begin{table}[H]
\centering
\small
\caption{Gate performance on the test sets. Each row is one test condition.
Columns: accuracy, F1, false-negative rate (fraction of fog frames routed as
clear) and false-positive rate (fraction of clear frames routed as fog).}
\label{tab:gate}
\begin{tabular}{@{}lcccc@{}}
\toprule
Test set & Accuracy & F1 & FNR (fog$\rightarrow$clear) & FPR (clear$\rightarrow$fog) \\
\midrule
Clear test set                  & 0.9141 & 0.0000 & 0.000 & 0.086 \\
Fog test set (CycleGAN)         & 1.0000 & 1.0000 & 0.000 & 0.000 \\
Fog test set (CycleDiff-75)     & 0.6612 & 0.7961 & 0.339 & 0.000 \\
\bottomrule
\end{tabular}
\end{table}

A 40-epoch comparison of routing strategies gives:
\begin{itemize}[leftmargin=1.4em]
  \item always-on U-Net: clear mAP@50--95 0.1592, fog 0.1510;
  \item gated U-Net: clear 0.1518, fog 0.1433;
  \item unfiltered baseline: clear 0.1464, fog 0.1234.
\end{itemize}
On this budget always-on filtering scores higher than gating. The comparison
is indicative rather than controlled, however: the always-on arm uses a YOLOv8s
detector and the gated arm a YOLOv12s detector, so detector capacity is
confounded with the routing strategy. The controlled comparison in this report
is the matched 100-epoch baseline-versus-filter table, which holds the detector
fixed. Gating remains an efficiency technique: it avoids the filter cost on
clear frames, but it does not improve accuracy and it can hurt when the gate
misclassifies.

\subsection{Ablations}
Training on clear data only and evaluating on fog confirms that the synthetic
weather is a real distribution shift: clear-only training reaches 0.1540 on
clear frames but only 0.0925 on CycleGAN fog and 0.1294 on CycleDiff-75 fog.
Training on fog only reaches 0.1313 (CycleGAN) and 0.1616 (CycleDiff-75) on fog.
The mixed pipelines are the best compromise.


%======================================================================
\section{FLIR ADAS v2: additional adverse conditions}
\label{sec:flir-adverse}

\subsection{Question}
The fog experiments above leave open whether the result is specific to fog, and
whether the detector is fragile to other thermal degradations. We therefore
screened eight further conditions on the same 3,749 matched frames, and
retrained only on conditions that actually hurt a clear-trained detector. This
is a different question from ``can a generator produce this weather?'': it asks
first whether the condition creates a measurable detection gap, and only then
whether adding it to training closes that gap.

\subsection{Conditions and protocol}
The conditions fall into two families. Fog (including mist and high-humidity
variants), rain, snow, thermal crossover and sensor degradation have a physical
approximation and could in principle be learned by a translation network. Small
targets, defocus and static obstruction alter geometry or sharp optical
properties, so they are generated directly: a translation network would
hallucinate structure rather than apply the required transform.
Figure~\ref{fig:07-conditions} shows one frame per condition.

\begin{figure}[H]
\centering
\includegraphics[width=\linewidth]{fig_07_conditions.png}
\caption{One FLIR test frame (000745) under the clear source and the eight
additional conditions at medium severity. Green boxes are the shared ground
truth. The conditions are, left to right and top to bottom: fog, rain, snow,
thermal crossover, sensor/dirty lens, small target, defocus and static
obstruction.}
\label{fig:07-conditions}
\end{figure}

The protocol is staged so that compute is not spent retraining on conditions the
detector already handles:
\begin{enumerate}[leftmargin=1.4em]
  \item Train \textbf{CM}, a YOLOv12s, on the 9,051 clear training frames
        (\textbf{C\_Tr}) with the same recipe used elsewhere in this report
        (40 epochs maximum, patience 10, batch 16, image size 640, seed 0).
  \item Evaluate CM on the clear test set \textbf{C\_T} and on each transformed
        test set \textbf{A\_T}. The screen statistic is the relative drop
        $(\mathrm{mAP}_{C_T} - \mathrm{mAP}_{A_T}) / \mathrm{mAP}_{C_T}$ in
        mAP@50--95.
  \item If the drop exceeds 5\,\%, generate the transformed training set
        \textbf{A\_Tr}, train \textbf{Mix\_M} on C\_Tr + A\_Tr (18,102 frames),
        and evaluate it on both C\_T and A\_T. The remaining gap is
        $\mathrm{Mix}(C_T) - \mathrm{Mix}(A_T)$, always clear minus adverse.
\end{enumerate}
Every evaluation uses exactly the same 3,749 frame indices and labels.
CM reaches mAP@50--95 $= 0.14966$ and mAP@50 $= 0.27546$ on C\_T.

\paragraph{Notation for Tables~\ref{tab:07-screen} and~\ref{tab:07-retrain}.}
\textbf{CM} (also written $C_M$) is the clear-trained YOLOv12s;
\textbf{Mix\_M} is the model trained on the mix of clear and
adversarial-condition images; \textbf{C\_T} is the 3,749-frame clear test set;
and \textbf{A\_T} is the same frames transformed into the adversarial condition.
\textbf{C\_Tr} and \textbf{A\_Tr} denote the corresponding training sets.

Table~\ref{tab:07-screen} reports the screen. Five conditions cross the 5\,\%
threshold: snow, sensor degradation, small targets, defocus and static
obstruction. Fog (2.14\,\%), rain (2.97\,\%) and thermal crossover (4.44\,\%)
fall below it and are not retrained under this protocol; they are recorded as
\emph{not selected}, not as harmless. The physical fog here is a
uniform-transmission veil and is milder than the CycleGAN fog of
Section~\ref{sec:flir}, so the below-threshold fog result is specific to this
implementation and severity.

\begin{table}[H]
\centering
\footnotesize
\setlength{\tabcolsep}{5pt}
\caption{Screen of the clear-trained YOLOv12s CM on the eight conditions
(3,749 frames each, mAP@50--95). The relative drop is
$(\mathrm{CM}_{C_T} - \mathrm{CM}_{A_T}) / \mathrm{CM}_{C_T}$. ``Selected''
means the condition is retrained in Table~\ref{tab:07-retrain}.}
\label{tab:07-screen}
\begin{tabular}{@{}lccccc@{}}
\toprule
Condition & CM C\_T & CM A\_T & Absolute drop & Relative drop & Selected \\
\midrule
Fog               & 0.14966 & 0.14645 & 0.00321 & 2.14\,\%  & no \\
Rain              & 0.14966 & 0.14521 & 0.00445 & 2.97\,\%  & no \\
Snow              & 0.14966 & 0.12652 & 0.02314 & 15.46\,\% & yes \\
Thermal crossover & 0.14966 & 0.14301 & 0.00665 & 4.44\,\%  & no \\
Sensor/dirty lens & 0.14966 & 0.09470 & 0.05496 & 36.72\,\% & yes \\
Small target      & 0.14966 & 0.12675 & 0.02291 & 15.31\,\% & yes \\
Defocus           & 0.14966 & 0.05373 & 0.09593 & 64.10\,\% & yes \\
Static obstruction& 0.14966 & 0.11775 & 0.03191 & 21.32\,\% & yes \\
\bottomrule
\end{tabular}
\end{table}

\subsection{Retraining the selected conditions}
Table~\ref{tab:07-retrain} reports the mixed models. Snow, sensor degradation,
defocus and static obstruction all reduce the gap substantially while keeping
clear accuracy close to CM. Snow and sensor degradation also improve clear
accuracy. Defocus and obstruction give up less than 0.003 clear mAP while
recovering 0.047 and 0.018 adverse mAP respectively.

\begin{table}[H]
\centering
\footnotesize
\setlength{\tabcolsep}{5pt}
\caption{Retrained models on the selected conditions (mAP@50--95, 3,749
frames). ``Mix gap'' is clear minus adverse. The small-target model produces a
\emph{negative} gap: it is better on A\_T than on C\_T.}
\label{tab:07-retrain}
\begin{tabular}{@{}lccccc@{}}
\toprule
Condition & Mix C\_T & Mix A\_T & Mix gap & CM gap & $\Delta$ clear \\
\midrule
Snow               & 0.16306 & 0.15806 & 0.0050  & 0.0231 & $+0.0134$ \\
Sensor/dirty lens  & 0.15057 & 0.12825 & 0.0223  & 0.0550 & $+0.0009$ \\
Small target       & 0.14323 & 0.37717 & $-0.2339$ & 0.0229 & $-0.0064$ \\
Defocus            & 0.14706 & 0.10070 & 0.0464  & 0.0959 & $-0.0026$ \\
Static obstruction & 0.14796 & 0.13521 & 0.0128  & 0.0319 & $-0.0017$ \\
\midrule
Fog, rain, thermal crossover & --- & --- & --- & --- & not selected \\
\bottomrule
\end{tabular}
\end{table}

\begin{figure}[H]
\centering
\includegraphics[width=\linewidth]{fig_07_pairs.png}
\caption{Source (top) and generated (bottom) frames for four selected
conditions. Sensor degradation adds fixed-pattern noise, drift, vignetting and
dirty-lens smudges; static obstruction is an alpha-blended fence/foliage mask;
small-target pixelation is applied only inside the labelled boxes; defocus is a
Gaussian point-spread convolution.}
\label{fig:07-pairs}
\end{figure}

\begin{figure}[H]
\centering
\includegraphics[width=\linewidth]{fig_07_results.png}
\caption{Left: relative mAP@50--95 drop of the clear-trained CM on each
condition; the dashed line is the 5\,\% screening threshold. Right: remaining
clear-minus-adverse gap for the retrained models (lower is better). Small
targets are the exception: the 8$\times$8 pixelation is a strong learnable cue,
so the mixed model is better on A\_T than on C\_T and the gap is negative.}
\label{fig:07-results}
\end{figure}

\subsection{The small-target result is a domain cue, not a label error}
The small-target mixed model reaches 0.3772 on A\_T but only 0.1432 on C\_T.
This is not a split leak: the clear and small-target test labels are identical
(3,749 files, 61,596 instances, the same nine classes). Reducing each target
crop to 8$\times$8 pixels and upsampling it creates blocky, high-contrast
targets that the detector can learn to exploit; on clear frames that cue is
absent, and clear mAP falls slightly ($-0.0064$). The result should therefore
be read as a warning about pixelation as a simulation, not as evidence that
small targets are easy.

\subsection{Incident and caveats}
The first sensor mixed-training attempt was killed by the scheduler
(SIGKILL, return code $-9$) after about two minutes while the dataloader was
caching 18,102 images in RAM under the 48\,GB allocation. It was retried with
on-the-fly decoding (no RAM cache) and 12 workers, and the completed run is the
one reported above. No condition is blocked.

These results use one seed and one severity (medium) per condition. The
conditions are simulations on 8-bit thermal JPEGs: FLIR has no real weather, no
depth and no radiometric temperature. Fog uses a uniform transmission by
default; crossover uses the local box-ring background as a proxy for a
sun-heated scene; obstruction is procedural unless real masks are supplied. A
severity sweep and a second seed are the obvious next steps.


%======================================================================
\section{Discussion}
\label{sec:discussion}

\paragraph{The measured weather gap depends on the protocol.}
On UNIRI-TID, mixed-weather training reduces the clear-to-fog APm gap to 0.094.
On FLIR, comparing 1,200 clear frames with 3,749 fog frames suggested a gap of
about 0.18, but matching the frames reduced it to 0.022. Any claim about weather
robustness must therefore state the evaluation protocol, and matched frames are
the minimum requirement.

\paragraph{Synthetic weather is learnable, but not all synthetic weather is equally useful.}
On M3FD, deterministic fog is essentially free after mixed training (gap
$-0.001$), whereas learned CycleGAN fog retains a 0.049 gap. The detector can
learn a fixed transform; it struggles with a transform whose texture and contrast
vary with the generator.

\paragraph{Only the learned filter helps, and only on the generator it was trained for.}
On FLIR, gamma, CLAHE and DCP do not improve detection, while the U-Net dehazer
does. The improvement is large enough to survive the stricter 100-epoch
comparison (+21.3\,\% fog mAP@50--95), and it also improves clear images. The
filter nevertheless fails to transfer to a different generator, and even
target-generator retraining does not beat the unfiltered baseline. This suggests
that reconstruction-based dehazing optimises the wrong objective for detection:
a visually cleaner image is not necessarily a more detectable one.

\paragraph{Gating is an efficiency technique.}
The gate is accurate on the generator it saw but misses a third of the fog frames
from a different generator. On the same 40-epoch budget, always-on filtering
scores higher than gated filtering, but the two arms use different detector
sizes (YOLOv8s versus YOLOv12s), so the comparison is indicative and not
controlled. The case for gating is therefore computational (the filter is
skipped on clear frames) rather than a demonstrated accuracy gain. We did not
record a latency benchmark, so the efficiency claim also remains qualitative in
this report.

%======================================================================
\section{Limitations and open work}
\label{sec:limitations}

\begin{itemize}[leftmargin=1.4em]
  \item \textbf{Synthetic fog.} FLIR has no real fog, so the FLIR conclusions are
        a proof of concept on generated weather, not a real-weather claim. UNIRI-TID
        is the only real-weather dataset, and its fog test set is small.
  \item \textbf{Single seed.} Every arm uses one random seed; no variance or
        confidence interval is reported.
  \item \textbf{Early stopping.} ``100 epochs'' means up to 100 with early
        stopping after 25 epochs without improvement. The effective budgets are
        55--65 epochs for the main arms.
  \item \textbf{Incomplete arm.} The filtered CycleDiff-75 pipeline stopped at
        epoch 27 of 100 and is not reported as a completed result.
  \item \textbf{Translation metrics.} PSNR, SSIM and LPIPS measure deviation from
        the clear source; no FID or human realism study was performed.
  \item \textbf{Invalid baseline.} The RT-DETR-l run on UNIRI-TID is invalid
        (29 epochs) and must not be quoted as a baseline.
  \item \textbf{Missing M3FD 3B data.} The Protocol-3B training images for M3FD
        are not shipped with the project; the result files are preserved, and the
        data can be regenerated from the split and synthetic-depth trees.
  \item \textbf{Real-time cost.} No latency or throughput measurements were
        recorded; the gating efficiency argument is therefore qualitative.
  \item \textbf{Additional-condition sweep.} The eight-condition extension uses
        one severity (medium) and one seed per condition. Fog, rain and thermal
        crossover fall below the 5\,\% screen at this severity; that is not
        evidence that the real weather is harmless. The physical fog is a
        uniform-transmission veil and is milder than the CycleGAN fog used in
        Section~\ref{sec:flir}.
  \item \textbf{Small-target cue.} Pixelating target crops to 8$\times$8
        pixels makes them blocky and more detectable, not merely smaller. The
        negative gap of the small-target mixed model is a property of this
        simulation, not a claim that small targets are easy.
\end{itemize}

%======================================================================
\section{Reproducibility}
\label{sec:repro}

\paragraph{Environment.}
Python 3.11.16, PyTorch 2.14.0 (CUDA 12.6), Ultralytics 8.4.138, NumPy 2.4.6,
OpenCV 5.0.0, scikit-learn 1.9.0, SciPy 1.17.1, PyYAML 6.0.3, timm 1.0.30 and
PyTorch Lightning 1.9.3. The full list is in
\texttt{environment/requirements.txt}. Training assumes one CUDA GPU; the
commands below must be run inside a Slurm \texttt{salloc} allocation as
\texttt{srun} steps, because direct \texttt{salloc} children were terminated by
the cluster after a few minutes.

\paragraph{Main commands.}
The exact entry points are:
\begin{verbatim}
# UNIRI-TID detector bench
cd experiments/05_detector_bench
python train_detector.py --model weights/yolov8s.pt \
    --name v8s --epochs 100 --cache ram
python eval/detect_eval.py --run v8s --split test

# M3FD
cd experiments/03_m3fd_weather
bash scripts/run_all_methods.sh

# FLIR, main comparison (inside an srun step)
cd experiments/02_flir_gate_ir
python -m src.build_direct init
bash scripts/run_main_100.sh
bash scripts/run_main_100_cyclediff.sh
\end{verbatim}

\paragraph{Artifacts.}
Each experiment directory contains its own \texttt{README.md},
\texttt{results/} (JSON metrics), \texttt{runs/} (weights and training curves) and
\texttt{logs/}. The consolidated result tables are in
\path{experiments/02_flir_gate_ir/results/REPORT.md},
\path{experiments/03_m3fd_weather/results/} and
\path{experiments/05_detector_bench/results/}. The figures in this report are
in \path{docs/report/figures/}.

\paragraph{Provenance.}
Every evaluation JSON stores the weights path and SHA-256 prefix, the dataset
YAML, the number of images, the split hash and the git state. The numbers in this
report were read from those files; no experiment was re-executed for the report.

%======================================================================
\section{Conclusion}

The project set out to test whether a learned enhancement stage makes thermal
detection robust to adverse weather, and whether gating that stage is worthwhile.
The evidence supports a nuanced answer. On real weather, mixed training already
closes most of the gap, so the headroom for a gate is small. On controlled
synthetic weather, a learned dehazer gives a clear and repeatable improvement on
the generator it was trained for (fog mAP@50 +17.5\,\%, gap $-43$\,\%), while
classical filters fail and cross-generator transfer fails. In an indicative
40-epoch comparison always-on filtering outscored gated inference, but the arms
differed in detector size, so that ranking is not controlled; the gate is
motivated by computation rather than by a demonstrated accuracy advantage. The most important methodological lesson is that weather robustness
must be measured on matched frames, because mismatched test sets can turn a
0.02 gap into an apparent 0.18 gap.

A second lesson is that not every simulated adverse condition is a robustness
problem. Of eight additional FLIR conditions, five crossed a 5\,\% screen; mixed
training then closed 52--78\,\% of the gap for snow, sensor degradation, defocus
and static obstruction. The exceptions are informative: a mild physical fog did
not hurt the detector, and small-target pixelation created a learnable cue that
produced a negative gap. Screening conditions before spending compute on
retraining is therefore both cheaper and more honest than assuming that every
simulated weather condition is harmful.

%======================================================================
\appendix
\section{Result and artifact map}
Table~\ref{tab:artifacts} lists the files from which the reported numbers and
figures were taken.

\begin{table}[H]
\centering
\scriptsize
\setlength{\tabcolsep}{3pt}
\caption{Where the numbers and figures in this report come from. Each row is one
artifact category; the right column lists the files. Paths are relative to the
project root.}
\label{tab:artifacts}
\begin{tabular}{@{}>{\raggedright\arraybackslash}p{0.28\linewidth}>{\raggedright\arraybackslash}p{0.66\linewidth}@{}}
\toprule
Content & Location \\
\midrule
UNIRI-TID detector results & \path{experiments/05_detector_bench/results/eval_*_test.json} \\
M3FD method results & \path{experiments/03_m3fd_weather/method1_classical/}, \path{method2_depth/}, \path{method3_cyclegan/} \\
M3FD detector results & \path{experiments/03_m3fd_weather/results/*.json} \\
FLIR filter screen & \path{experiments/02_flir_gate_ir/results/*_60_*.json} \\
FLIR main 100-epoch results & \path{experiments/02_flir_gate_ir/results/*_100_*_vtt.json} \\
FLIR consolidated tables & \path{experiments/02_flir_gate_ir/results/REPORT.md} \\
Gate metrics & \path{experiments/02_flir_gate_ir/results/gate_cyclegan.json} \\
Translation metrics & \path{experiments/01_flir_weather_translation/method4_cyclediff/logs/} \\
Additional FLIR condition screen & \path{experiments/07_flir_adverse_conditions/results/eval_CM_*.json} \\
Additional FLIR retraining (evals) & \path{experiments/07_flir_adverse_conditions/results/eval_Mix_*.json} \\
Additional FLIR retraining (runs) & \path{experiments/07_flir_adverse_conditions/runs/Mix_*_yolo12s/} \\
Additional FLIR condition report & \path{experiments/07_flir_adverse_conditions/results/FINAL_REPORT.md} \\
Additional condition generation & \path{experiments/07_flir_adverse_conditions/data/generated/*/medium/generation.json} \\
Report figures & \path{docs/report/figures/} \\
\bottomrule
\end{tabular}
\end{table}

\begin{thebibliography}{9}
\bibitem{uniri} K.~Kristo et al., ``Thermal imaging dataset for person detection''
(UNIRI-TID), 2020.
\bibitem{m3fd} J.~Liu et al., ``Target-aware dual adversarial learning and a
multi-scenario multi-modality benchmark to fuse infrared and visible for object
detection'' (M3FD), CVPR 2022.
\bibitem{flir} Teledyne FLIR, ``FLIR ADAS v2 thermal dataset,'' public release.
\bibitem{cyclegan} J.-Y.~Zhu, T.~Park, P.~Isola and A.~A.~Efros, ``Unpaired
image-to-image translation using cycle-consistent adversarial networks,''
ICCV 2017.
\bibitem{cyclediff} CycleDiff: cycle-consistent diffusion for unpaired
image-to-image translation, arXiv:2508.06625.
\bibitem{clahe} K.~Zuiderveld, ``Contrast limited adaptive histogram
equalization,'' Graphics Gems IV, 1994.
\bibitem{dcp} K.~He, J.~Sun and X.~Tang, ``Single image haze removal using
dark channel prior,'' CVPR 2009.
\bibitem{yolo} Ultralytics, YOLOv8/YOLOv12 documentation and implementation.
\end{thebibliography}

\end{document}
